\section{Emergent relativity and effective Lorentz symmetry}
\label{sec:emergent-relativity}

The continuum model is written with an external time parameter $t$, so Lorentz invariance is not postulated
microscopically. Nevertheless, in the \emph{linear wave regime} an effective Lorentz symmetry can arise as a
kinematic symmetry of the dominant dispersion branch that controls information transport in a given regime.

\subsection{Linear-wave cone and an effective light speed}

Linearizing \Cref{eq:eom} about a homogeneous background yields a multi-branch wave system.
For isotropic elastic media, one expects (at long wavelengths) transverse and longitudinal branches.
In the StVK baseline around a homogeneous state, the characteristic speeds take the familiar form
\begin{equation}
  c_T^2 \sim \frac{\mu}{\rho_m},
  \qquad
  c_L^2 \sim \frac{\lambda+2\mu}{\rho_m},
  \label{eq:cl-ct}
\end{equation}
up to the details of the chosen background and the pre-stress encoded by $\alpha$.
More generally, for each polarization branch $b$ one has the long-wavelength dispersion
\begin{equation}
  \omega_b^2(\veck)\;\approx\; c_b^2\,|\veck|^2 \quad \text{for } |\veck|\to 0,
  \label{eq:linear-dispersion}
\end{equation}
which defines an effective wave cone for that branch.

In regimes where a single branch dominates transport and dispersion is negligible,
the wave equation is approximately invariant under Lorentz transformations built from the corresponding
speed $c_b$. In particular, for the dominant branch (denote its speed by $c$) one obtains the familiar
kinematic relations for boosted frames,
\begin{align}
t' &= \gamma\left(t-\frac{vx}{c^2}\right),&
x' &= \gamma(x-vt),&
\gamma &= (1-v^2/c^2)^{-1/2},
\label{eq:lorentz}
\end{align}
as an effective symmetry of the linear propagation sector, not as a microscopic postulate.

\subsection{Analogue-metric viewpoint}

An equivalent viewpoint, common in analogue-gravity systems, is that linear excitations experience an
effective metric structure with interval
\begin{equation}
  ds^2 = -c^2 dt^2 + g^{(\text{eff})}_{ij}(x,t)\,dx^i\,dx^j,
  \label{eq:effective-metric-interval}
\end{equation}
where $g^{(\text{eff})}_{ij}$ depends on the background configuration about which one linearizes.
In the present model there are two natural geometric objects: the induced brane metric $g_{ij}(\vecX)$
and the branch-dependent effective coefficients of the linearized operator.
In particular, slowly varying transverse profiles $u=X^4$ modify intrinsic distances already at the level of the induced metric (\Cref{eq:metric-near-identity,subsec:gravity-channel}), and can therefore be read as a background that focuses or defocuses wave transport in close analogy to a weak gravitational field.
Both encode how background deformations can bend or modulate wave propagation, in close analogy to
analogue-gravity constructions \citep{Barcelo2011AnalogueGravity,Volovik2003HeliumDroplet}.

\subsection{Regime of validity and expected deviations}

Effective Lorentz symmetry is expected to break when:
(i) multiple branches contribute comparably (mode mixing),
(ii) dispersion becomes significant at shorter wavelengths, or
(iii) geometric/material nonlinearity becomes important at larger amplitudes.
In the present theory these deviations are not treated as flaws; they are structural consequences of a
preferred-time substrate and therefore delimit the domain where a relativistic effective description is expected
to apply.
